\section{Conclusions and Discussions}
We introduce Language-Action Pre-training (LAP), a scalable VLA pre-training recipe that predicts actions in natural language, yielding representations that transfer effectively across robot embodiments. We further present \textsc{LAP-3B}, which achieves substantial zero-shot generalization on real robots—delivering roughly a 30\% absolute improvement over prior action representations and, to our knowledge, the first substantial zero-shot success rate on novel embodiments without fine-tuning. \textsc{LAP-3B} also exhibits stable training, favorable scaling, and strong fine-tuning efficiency.


\noindent\textbf{Broader implications.}
LAP opens the possibility of training VLMs \emph{explicitly for robotics}, enabling robot (and potentially non-robot) interaction data to be seamlessly mixed with standard VLM corpora to learn models that reason, communicate, and act within a unified representation space.

\noindent\textbf{Limitations and future work.}
Despite these promising results, several limitations remain. While LAP can be applicable to substantially more heterogeneous robot systems and data sources, this paper focuses only on zero-shot transfer across single-arm manipulators. For example, LAP can be easily adapted to bimanual robots, which actions can be represented through coordinated dual-arm descriptions (e.g., Left Arm: move forward 5 cm; Right Arm: rotate clockwise 20 degrees). Another key advantage of language-actions is that they naturally operate at varying levels of precision and are more
tolerant to less accurate action labels, enabling the use of data such as human pose tracking, UMI, or internet video where precise low-level control signals are difficult to obtain. Systematically investigating these directions at scale remains an important avenue for future work.

Second, although \textsc{LAP-3B} adapts efficiently to moderately dexterous tasks such as \emph{Fold Towel and Put in Basket} and \emph{Hang Tape on Rack}, we have not yet evaluated regimes requiring substantially higher control frequency or extreme precision (e.g., fast reactive control, or fine-grained deformable object manipulation). Extending language-action supervision to these settings—potentially through hierarchical or multi-scale action representations—remains an open challenge.